\section{Testable Propositions and Experimental/Data Strategies}

\subsection{Quasi-Periodic Spectrum and Golden Sector Discrimination}

\textbf{Proposition P1 (Quasi-periodic spectral lines).} If the three-dimensional observable set derives from an irrational-orientation cut-and-project model set, its diffraction spectrum or structure factor should exhibit quasi-periodic characteristics dominated by discrete pure-point spectrum, and should manifest as predictable ``phase slip'' upon variation of complementary phase offset $u$ rather than random drift.\\
\textbf{Proposition P2 (Statistical signal of golden sector preference).} Under fixed resource constraints, if the scan slope is in the minimal-resonance sector, then the differential uniformity index of prefix readout sequences (such as star discrepancy or gap distribution) should achieve comparable extreme performance within the same complexity class, thus enabling statistical discrimination from other irrational slopes.

\subsection{Degeneracy Fingerprint of Stable Sector Folding}

\textbf{Proposition P3 (Folding degeneracy histogram).} Under given window length $m$ and fixed folding rule, the preimage cardinality distribution (degeneracy histogram) of $\Fold_m$ is discrete and reproducible combinatorial invariant; different experimental platforms, as long as they share the same protocol syntax constraint, should exhibit degeneracy fingerprints transferable across systems. This proposition can be verified by window-based statistical reconstruction of discrete readout records.

\subsection{Verification of Delayed-Choice Equivalence Class Reallocation}

\textbf{Proposition P4 (Equivalence class reallocation).} If delayed change of measurement settings corresponds to changing readout kernel/projection parameters, then the change in observational statistics should be explainable by ``reallocation of equivalence class decomposition,'' and should satisfy computable conservation relations determined by the folding rule (such as total weight conservation, accessible migration graph constraints between specific types). This proposition requires type-level re-encoding of experimental readout, rather than comparison merely at wave-packet image level.

\subsection{Closed-Loop Stability of Adaptive Observer}

\textbf{Proposition P5 (Closed-loop convergence domain).} If the observer update law approximately executes gradient descent of the mismatch functional, then there exist observable convergence domains and phase transition boundaries: when external disturbance or resource constraints cross thresholds, the system transitions from stable-type maintenance region to defect accumulation region, and shows measurable mismatch growth and type collapse. This can be verified by implementing parameter feedback in controllable systems (such as quantum optical feedback, programmable media, or discrete cellular automaton experiments) and statistically measuring stable sector maintenance rates.
